\chapter*{Abstract}

This thesis presents \emph{Streamify}, a project built to learn, hands-on, how a modern streaming platform works and to apply distributed-systems concepts to a real system, rather than studying them in the abstract. A continuous stream of synthetic music-streaming events is carried end to end: ingested with ordering and fault-tolerance guarantees, processed continuously as it arrives rather than on a batch schedule, landed into a storage layer that supports both transactional updates and analytical queries, transformed into an analytics-ready model, and visualized and monitored throughout -- all under a small, fixed compute budget rather than assumed elastic scale.

\noindent Building this meant confronting, and resolving, the concrete distributed-systems problems that appear once a platform has more than one producer and consumer of the same data: partitioning and replication for ordered, fault-tolerant ingestion; checkpointed, exactly-once processing so a crash never silently loses or duplicates an event; a storage layer that serves frequent small writes and large analytical reads without one starving the other; and observability that shows whether the system is keeping up with its input, not just whether it is running.

\noindent Every major technology choice was made deliberately against concrete alternatives and the constraints of a single-machine deployment, rather than assumed as a given (Chapter~\ref{ch:related-work-background}). This thesis documents that reasoning alongside the resulting architecture (Chapter~\ref{ch:proposed-solution}) and what building and operating the system revealed about how these distributed-systems concepts behave in practice (Chapter~\ref{ch:conclusion}).
